\begingroup
\fontsize{11.4}{14}\selectfont
\noindent Deribit, 1 мая 2022~г. --- 14 ноября 2026~г.\par
\begin{center}\small\begin{tabular}{lrrrr}\toprule
Стратегия & 2024 & 2025 & 2026 & CAGR\\\midrule
Стрэддл + грид & 0,00 & +3,14 & +0,18 & +1,15\\
Грид, ежегодный подбор & $-0{,}12$ & $-0{,}06$ & +0,19 & +0,0016\\
Стрэддл + дельта-хедж & 0,00 & +2,18 & +0,11 & +0,79\\
Тренд за 42 дня & +29,98 & $-6{,}75$ & +21,40 & +14,40\\
Тренд + стрэддл с гридом & +14,99 & $-2{,}48$ & +11,73 & +8,17\\\bottomrule
\end{tabular}\end{center}
\smallskip
\textbf{Гипотезы и результаты.} В таблице доходность в процентах после модельных издержек. Каждый счёт стартует с 150\,000 USDC 1 января 2024 г. Годовые столбцы отсчитываются от начала года; $\mathrm{CAGR}=(E_T/E_0)^{365/d}-1$ за $d$ суток.
\begin{enumerate}
\setlength{\itemsep}{4pt}\setlength{\parsep}{0pt}\setlength{\topsep}{3pt}
\item \textbf{Покупка стрэддла.} Гипотеза: $\mathbb E[|D_T-K|\mid\text{сигнал}]>P+C$, где $P$ и $C$ --- премия и расходы на пару. На шести исполненных входах чистый стрэддл дал $+1{,}24\%$ годовых. У связки с гридом один цикл дал 3\,525 из 4\,984 USDC прибыли; без него осталось $+0{,}34\%$ годовых.
\item \textbf{Сетка.} Если возвраты к центру оплачивают инвентарь и издержки, то на общих входах $\Pi_{\text{стрэддл+грид}}-\Pi_{\text{стрэддл}}>0$. В исходном опыте задания 3 эта разность равна $-7\,763$ USDC за 36 циклов. У самостоятельного итогового грида 35 циклов дали всего $+6{,}90$ USDC. Без лучшего цикла CAGR $-0{,}0302\%$; при участии 10\% вместо 25\% и проходе и проскальзывании по 5 б.п. CAGR $-0{,}054\%$.
\item \textbf{Частичный хедж.} Гипотеза: при $H^*=-0{,}5Q\Delta$ просадка меньше, чем у чистой пары. На тех же шести входах она снизилась с $0{,}96\%$ до $0{,}68\%$ по дневным отметкам и выходам. Цена снижения риска --- 1\,953 USDC недополученной чистой прибыли. Все три бота задания с 2024 г. дали менее 4\% годовых.
\item \textbf{Тренд.} Гипотеза продолжения движения: $\mathbb E[w_t r_{t+1}]>0$, где знак $w_t$ задаёт доходность за 42 дня. С мая 2022 г. чистый CAGR составил $+16{,}68\%$ при комиссии 3,5 б.п. и $+15{,}97\%$ при 5 б.п. Без лучшей серии при исходной комиссии осталось $+9{,}65\%$; максимальная пятиминутная просадка составила $42{,}60\%$ на полной серии.
\item \textbf{Портфель 50/50.} Гипотеза: опционная часть повышает итог относительно тех же 75\,000 USDC в тренде и 75\,000 USDC в деньгах под 0\%. С 2024 г. капитал портфеля составил 187\,946,98 USDC против 185\,366,18 USDC у такого контроля. Вклад опционной части $+2\,580{,}79$ USDC; без лучшей серии одной из частей CAGR портфеля $+2{,}46\%$.
\end{enumerate}
\sub{Следующая проверка} В исходном опыте за первый час после сигнала покупки обеих ног видны в 17 из 36 циклов. Лишь в двух по обеим ногам набрался объём в пределах лимита. Неизвестно, получили бы наши заявки исполнение в очереди. Зафиксируем календарь и сигнал; на новых входах запишем синхронные bid/ask, глубину, очередь, задержку и время с одной открытой ногой. Проверим падение IV до расчёта, скачки базиса BTC/USD--USDC, расхождение с TWAP у страйка, спред при ролле и доступность маржи при сбое биржи. Критерий продолжения: положительный чистый вклад грида против чистой пары на новых общих входах и прибыль связки при худшем исполнении. При отдельном старте в январе 2025 г. исходные автономные связка, грид и частичный хедж дали $-4{,}47\%$, $-2{,}09\%$ и $-5{,}96\%$ годовых. Для вывода о 16\% у связки нужна новая серия сделок.
\par\smallskip{\small Данные и спецификации: \href{https://docs.deribit.com/api-reference/market-data/public-get_last_trades_by_currency_and_time}{Deribit, история сделок}; \href{https://support.deribit.com/hc/en-us/articles/31424939096093-Inverse-Options}{инверсные опционы}; \href{https://support.deribit.com/hc/en-us/articles/25944746248989-Fees}{тарифы}; \href{https://docs.tardis.dev/historical-data-details/deribit}{Tardis, котировки}.}
\endgroup
\newpage

\begingroup
\fontsize{10.25}{11.55}\selectfont\setlength{\leftmargini}{1.55em}
\setlength{\parskip}{1pt}\setlength{\abovedisplayskip}{2pt}\setlength{\belowdisplayskip}{2pt}
\textbf{\large Правила итоговых ботов и статистика входов}\\[2pt]01.01.2024--14.11.2026, UTC.\par
$E$ --- капитал в USDC, $S$ --- индекс BTC, $K$ --- страйк, $D_T$ --- цена расчёта.
$[x]_a=a\lfloor x/a\rfloor$. Precision $=TP/(TP+FP)$, recall $=TP/(TP+FN)$.
Одно исполнение пакета STRD считается одной сделкой; экспирация --- денежный расчёт. Метрики сигнала STRD общие для пп. 1, 2, 4 и опционной части п. 6; метрики тренда --- для пп. 5 и 6.
\begin{enumerate}
\setlength{\itemsep}{2pt}\setlength{\parsep}{1pt}\setlength{\topsep}{2pt}
\item \textbf{Чистый стрэддл.} Каждый час при свободном счёте выбирается пакет call + put одного $K,T$: $|K/S-1|\le3\%$, до расчёта от 6 часов до 7 суток. Берётся максимум объёма за 24 часа среди пакетов с доступной IV; при равенстве --- ближайший страйк, затем ближайшая экспирация. $M=C_{BS}+P_{BS}$ по IV каждой ноги при нулевой ставке. Для $j=1,\dots,13$ недель и уже завершённых аналогов
\[
 X_j=|S_t D_{T-7j}/S_{t-7j}-K|,\;
 \bar X=\frac{\sum_j X_j-\max_j X_j}{n-1},\ n\ge8,\;
 B=\big[\min(0{,}975M,0{,}90\bar X)/S_t\big]_{0{,}0001}.
\]
Для ноги $b=\pm1$ (call/put) $d_1=\ln(S/K)/(\sigma\sqrt\tau)+\sigma\sqrt\tau/2$, $d_2=d_1-\sigma\sqrt\tau$, цена $b[S\Phi(bd_1)-K\Phi(bd_2)]$, $\Delta=b\Phi(bd_1)$, $\Gamma=\phi(d_1)/(S\sigma\sqrt\tau)$, $\tau$ в годах. С 1-й минуты до конца часа покупка по $B$ при первой встречной продаже $p\le B-0{,}0001$ BTC с доступным лотом.
$Q=[\min(E/S_t,E/S_{\rm fill},0{,}08E/c,0{,}25V)]_{0{,}1}$ BTC, где $V$ --- объём принта. Стоимость пары $c=1{,}001\sum_i(a_i+\min(0{,}0003S_{\rm fill},0{,}125a_i))$, $a_i=S_{\rm fill}(p_i+(B-p)/2)$ для ног $i=C,P$. Удержание до выплаты $Q|D_T-K|$. \textbf{6 сделок.}
На 99 встречных продажах от 0,4 BTC сигнал --- покупка, событие --- выплата выше цены принта: $|D_T-K|>pS_{\rm fill}$. Precision $5/6=83{,}33\%$, recall $5/32=15{,}625\%$.
\item \textbf{Стрэддл + грид.} Покупка и расчёт по п. 1. С первой полной минуты после покупки цена-якорь $F$ берётся по последнему закрытию, затем каждые 5 минут. $h=F\max(0{,}0015;\,0{,}35\sqrt{\sum_{i=1}^{12}r_{i,5}^2})$, где $r$ --- лог-доходности перпетуала.
\par Объём уровня $u=[\min(0{,}10E/(3F),0{,}75Q/3,0{,}5Q\Gamma h)]_{0{,}001}$, $\Gamma=\Gamma_C+\Gamma_P$. При позиции $nu$ центр $A=F-0{,}5nh$; покупка по $\lfloor A-h\rfloor$, продажа по $\lceil A+h\rceil$, $|n|\le3$. Нужны проход 2 б.п. и объём свечи $\ge4u$; двусторонний проход пропускается. За 30 минут до $T$ начинается закрытие инвентаря. \textbf{369 сделок}: 6 пакетов + 363 исполнения перпетуала.
\item \textbf{Самостоятельный грид.} С 1-й минуты часа $F=P_{\rm inv}\,\mathrm{EWMA}_{1/60}(P_{\rm lin}/P_{\rm inv})$ по завершённым минутам с торговлей обоими контрактами; $P$ --- предыдущая цена линейного перпетуала. Уровень $u=[0{,}10E/(2F_0)]_{0{,}001}$, максимум $2u$. В 2024 г. при свежей сделке и $P/F-1<-0{,}0025$ покупка по $P(1-0{,}0015)$; при $P/F-1>0{,}0025$ зеркальная продажа. Добор ещё через 15 б.п. от первого исполнения при сохранении расхождения. В 2025--2026 гг. заявки $F(1\pm0{,}012)$, добор $F(1\pm0{,}024)$. Выход к $F$ или закрытие с 20-й минуты (2024--2025), 50-й (2026). Лимиты наружу до 1 USDC, проход 2 б.п., участие до 25\% встречного объёма. \textbf{71 сделка, 35 циклов.} Сигнал --- первая покупка в цикле, событие --- последняя цена исполнения выше первой до проскальзывания. Precision $9/17=52{,}94\%$, recall $9/18=50\%$.
\item \textbf{Стрэддл + дельта-хедж.} Пакет по п. 1. Каждую минуту цель $H^*=-0{,}5Q(\Delta_C+\Delta_P)$, полоса $H^*\pm0{,}15Q$. Ниже полосы покупка до нижней границы, выше --- продажа до верхней; округление внутрь до 0,001 BTC, исполнение со следующей минуты. С $T-30$ минут цель ноль. \textbf{280 сделок}: 6 пакетов + 274 исполнения перпетуала.
\item \textbf{Тренд за 42 дня.} В 00:00 $r_t=\ln(P_t/P_{t-1})$, $v_{k,t}=2^{-1/k}v_{k,t-1}+(1-2^{-1/k})r_t^2$, $k\in\{7,30\}$; $s=\max(0{,}20,\sqrt{365(0{,}65v_7+0{,}35v_{30})})$.
$w=\min(1,0{,}45/s)\,\mathrm{clip}(\ln(P_t/P_{t-42})/(s\sqrt{42/365}),-1,1)$.
С 00:05 на первом доступном открытии $P_o$ цель $H=\mathrm{sign}(w)[|w|E/P_o]_{0{,}0001}$. Покупка при росте цели относительно позиции, продажа при снижении, закрытие в конце теста. \textbf{1036 сделок.} На 1048 дневных сигналах прогноз роста $w>0$; факт --- следующая цена исполнения дневной цели выше текущей (для последнего дня --- конечное закрытие). Precision $305/591=51{,}61\%$, recall $305/528=57{,}77\%$.
\item \textbf{Тренд + стрэддл с гридом.} $E_{2,0}=E_{5,0}=75\,000$, $E_t=E_{2,t}+E_{5,t}$. Сделки по пп. 2 и 5 от капитала каждой части, без переводов. \textbf{1466 сделок}: 6 пакетов, 433 исполнения сетки, 1027 трендовых.
\end{enumerate}
\textbf{Исходный протокол для сравнения на стр.\,3.} Каждую вторую пятницу с 06.05.2022 в 12:00 UTC: ближайший общий страйк в пределах 3\% от $S$, расчёт через 14 дней в 08:00. Нужны сделки обеих ног за 6 часов. Если размах восьми дневных отметок $S$ за неделю не меньше $P_0=\sum_i[M_i+\max(0{,}025M_i,0{,}00005S)]$, в 12:01 покупаем $Q=[\min(0{,}08E/D_0,E/S)]_{0{,}1}$ пар до расчёта ($D_0$ --- премия $P$ со входными расходами). Шаг $h=0{,}005S$ до 0,1 USDC; границы $K\pm Nh$, $N=\lfloor P/h\rfloor$. Уровень $q=[Q\min(0{,}5\Gamma h,0{,}75/(N+|i_0|))]_{0{,}0001}$, $i_0$ --- начальное смещение в шагах. На соседнем нижнем уровне покупка, на верхнем продажа. Пробой по минутному закрытию: выход на следующем открытии; иначе выход с 07:30 перед расчётом. У бенчмарка полной дельты цель $-Q\Delta$, порог $0{,}05Q$. Автономный грид входил при $|\sum_1^7 r_j|/\sum_1^7|r_j|\le1/3$ и держал остаток до расчётного окна; автономный хедж возвращался к $-0{,}5Q\Delta$ при отклонении $0{,}15Q$.
\endgroup
